\documentclass[11pt]{article}

\usepackage{amsmath,amssymb}
\usepackage{algorithm}
\usepackage[noend]{algpseudocode}
\usepackage{float}
\usepackage{graphicx}
\usepackage{geometry}
\usepackage{xurl}
\usepackage{booktabs}
\geometry{margin=1in}

\title{Construction of SPF Datasets for the ME3AI Project}
\date{}
\newcommand{\worktreerootrelative}{../..}
\newcommand{\forecastrevisionoutputroot}{\worktreerootrelative/data-pipeline/output/forecast_revision}

\begin{document}
\maketitle

\begin{abstract}
This document specifies the procedure for constructing Survey of Professional
Forecasters (SPF) datasets used in this project. The SPF is a quarterly survey
of macroeconomic forecasts conducted by the Federal Reserve Bank of
Philadelphia; data are provided as Excel workbooks (median and mean levels and
growth rates, probabilities, dispersion, and microdata). We describe how to
obtain the official releases, which variables and vintages to use, how to
validate and normalize the data into the project's cleaned third normal form
(3NF) structure, and how to link them to the pipeline input and output folders.
This document serves as the single source of truth for SPF dataset construction
and reproducibility.
\end{abstract}

\begin{table}[h]
\centering
\caption{Notation and terms used in this document.}
\label{tab:notation}
\begin{tabular}{@{}ll@{}}
\toprule
Symbol or term & Definition \\
\midrule
SPF & Survey of Professional Forecasters \\
3NF & Third normal form (database normalization) \\
BEA & Bureau of Economic Analysis \\
NIPA & National income and product accounts \\
     & (BEA advance report) \\
\texttt{var\_path} & URL path segment for a variable (lowercase; \\
     & e.g., \texttt{ngdp}, \texttt{cpi10}) \\
VAR & Placeholder for variable code in filenames (e.g., NGDP, CPI10) \\
\texttt{overwrite} & Policy: \texttt{overwrite} or \texttt{skip-if-exists} \\
\bottomrule
\end{tabular}
\end{table}

\section{Data source and scope}

The Survey of Professional Forecasters (SPF) is a quarterly survey of
professional forecasters, formerly run by the American Statistical Association
and the National Bureau of Economic Research (from 1968:Q4) and taken over by
the Federal Reserve Bank of Philadelphia in 1990:Q2. The Philadelphia Fed
publishes documentation and data files; the canonical reference for variable
definitions, file layout, and timing is the \emph{SPF Documentation}
(Philadelphia Fed, latest update March 2024), held in this project as
\texttt{reference/spf-documentation.pdf}.

The survey is timed to the release of the Bureau of Economic Analysis (BEA)
advance report of the national income and product accounts (NIPA):
questionnaires are sent after that report; deadlines fall in the second to
third week of the middle month of each quarter; results are released before
the BEA second report. For the construction of our datasets, we use the
survey date (year and quarter) as the primary time index. The first survey
fully conducted in real time by the Philadelphia Fed is 1990:Q3. Earlier
vintages (1968:Q4 to 1990:Q2) may be used with the caveats noted in the
official documentation.

\section{Published files and variables}

The Philadelphia Fed provides SPF data \emph{by variable}: each variable has a
dedicated data page from which Excel (and other) files can be downloaded. The
variable index is at \url{https://www.philadelphiafed.org/surveys-and-data/data-files};
clicking a variable (e.g., NGDP, CPI10, RGDP, UNEMP) leads to a page at
\[
\texttt{https://www.philadelphiafed.org/surveys-and-data/}\textit{var\_path},
\]
where $\textit{var\_path}$ is the variable identifier in lowercase (e.g.,
\texttt{ngdp}, \texttt{cpi10}, \texttt{rgdp}). The sole exception is CPI
inflation, which uses \texttt{cpi-spf}. Variable names and paths are listed on
the data-files index (e.g., ``10-Year CPI Inflation Rate (CPI10)'' maps to
\texttt{/cpi10}).

On each variable page, the following file types are typically offered
(availability may vary by variable):

\begin{itemize}
  \item \textbf{Documentation}: \\
    Link to the SPF documentation PDF (same for all variables).
  \item \textbf{Measures of Cross-Sectional Forecast Dispersion}: \\
    \texttt{Dispersion\_}\textit{VAR}\texttt{.xlsx} (e.g.,
    \texttt{Dispersion\_CPI10.xlsx}, \texttt{Dispersion\_NGDP.xlsx}).
  \item \textbf{Median Responses}: \\
    \texttt{Median\_}\textit{VAR}\texttt{\_Level.xlsx} (e.g.,
    \texttt{Median\_NGDP\_Level.xlsx}, \texttt{Median\_CPI10\_Level.xlsx}).
  \item \textbf{Mean Responses}: \\
    \texttt{Mean\_}\textit{VAR}\texttt{\_Level.xlsx}.
  \item \textbf{Annualized Percent Change of Median/Mean Responses}: \\
    \texttt{Median\_}\textit{VAR}\texttt{\_Growth.xlsx} and
    \texttt{Mean\_}\textit{VAR}\texttt{\_Growth.xlsx}. For variables that have
    growth rates (e.g., NGDP, RGDP); not all variables such as CPI10 offer
    growth files.
  \item \textbf{Individual Responses}: \\
    \texttt{Individual\_}\textit{VAR}\texttt{.xlsx}.
\end{itemize}

Filenames use the variable code as on the website (e.g., NGDP, CPI10). The
automation must fetch the variable page, extract the download links for the
desired file types, and save each file to \texttt{data-pipeline/input/} under
that same filename (no timestamps). Variable names, horizons, and caveats are
defined in the SPF documentation; we record the exact variable names and
definitions in a project codebook.

\section{Download automation: inputs and outputs}

We automate retrieval of SPF files from the Philadelphia Fed website so that
raw files are saved to \texttt{data-pipeline/input/} without manual download.
Given a variable name (e.g., CPI10, NGDP), the data are hosted on that
variable's page at \url{https://www.philadelphiafed.org/surveys-and-data/}%
$\textit{var\_path}$.
The automation fetches the page, parses the HTML for the desired download
links (Dispersion, Median Level, Mean Level, Median Growth, Mean Growth,
Individual), and downloads each file to the input folder under the same
filename as on the site (e.g., \texttt{Median\_CPI10\_Level.xlsx}). When
parsing links, the implementation must match \texttt{<a>} tags where
\texttt{href} may appear after other attributes (e.g.\ \texttt{class}), since
the Philadelphia Fed markup may use that order.

\paragraph{Inputs.}
A list of variable names (e.g., CPI10, RGDP, UNEMP, NGDP) as on the
Philadelphia Fed data-files index. Optional inputs: which file types to
download (e.g., median level, mean level, dispersion, individual); overwrite
policy $\texttt{overwrite} \in \{\texttt{overwrite}, \texttt{skip-if-exists}\}$
(whether to overwrite an existing file or skip); output directory (default
\texttt{data-pipeline/input/}).

\paragraph{Outputs.}
Each selected file is saved under \texttt{data-pipeline/input/} with the same
filename as on the website (no timestamps or version suffixes). Each file is
downloaded at most once per variable; if multiple variables were to share the
same filename, the implementation should deduplicate so that the file is
written once.

Three logical functions suffice: (1) given a variable name, return the variable
page URL (lowercase path, with exception CPI $\to$ \texttt{cpi-spf}); (2) given
a variable page URL and requested file types, fetch the page and return a list
of (filename, download URL) pairs; (3) given a download URL and local path,
download the file subject to the overwrite policy. The top-level procedure
loops over the requested variables, obtains the link list for each, and
downloads each file to \texttt{data-pipeline/input/}.

\section{Pipeline and 3NF}

Our data pipeline (see \texttt{docs/workflow/data\_pipeline.md}) has three
stages: \textbf{input} (raw data), \textbf{cleaned} (third normal form, 3NF),
and \textbf{output} (analysis results). SPF data enter as raw files in the
input folder and are transformed into cleaned tables that satisfy 3NF.

\paragraph{Input.}
Place the official Excel workbooks in the pipeline input folder. Do not
modify the original files; preserve them as the source of truth.

\paragraph{Cleaned.}
The current implementation focuses on \textbf{individual (microdata)} forecasts
only. Cleaned output is CSV: \texttt{forecast\_individual.csv} in \textbf{wide}
form (primary key: survey\_year, survey\_quarter, forecaster\_id; one
column per forecast horizon, all horizons in the same row) and \texttt{forecaster\_survey.csv}
(primary key: survey\_year, survey\_quarter, forecaster\_id; \texttt{industry}).
No separate survey table is used; the set of (survey\_year, survey\_quarter) is
derivable from \texttt{forecast\_individual}. Level/mean/median and probability
tables may be added in a later stage.

\paragraph{Output.}
Analyses read from the cleaned folder and write results to the output
folder, organized by analysis type.

Before treating a constructed SPF dataset as final: verify that raw files match
the expected Philadelphia Fed release; confirm variable names and column
layouts align with the SPF documentation; after transformation, check row
counts and primary keys; record data source, retrieval date, and any
filtering in the codebook or a README in the cleaned folder.

\section{Prepare dataset for forecast revision regression}

In this section, all constructed objects are indexed by a forecast unit $u$ and
a survey date $s=(y,q)$. Here the unit index $u$ corresponds to
\texttt{forecaster\_id}.
Let $s^{-}$ denote the previous survey date corresponding to $s$: if
$q \in \{2,3,4\}$ then $s^{-}=(y,q-1)$, while if $q=1$ then
$s^{-}=(y-1,4)$. With this notation, the three key constructed objects are the
long-term inflation expectation $x_{u,s}$, inflation news $n_{u,s}$, and the
reputation measure $\rho_{u,s}$.

\subsection{Construct long-term inflation expectation}

Following Appendix A.1.3 in Aruoba, let $\pi_{a \to b}$ denote annualized
quarterly inflation between quarters $a$ and $b$, and use the same
continuous-compounding approximation. The SPF 10-year CPI forecast is a
\emph{fixed-event} forecast: in each survey quarter it refers to inflation over
10 years starting from the previous year's Q4. Hence part of the reported
object is already realized by the survey date and must be removed before using
it as a forward-looking measurement equation.

In the cleaned wide table \texttt{forecast\_individual.csv}, both the raw
10-year expectation and the realized short-horizon inflation term are taken from
the same survey row: column \texttt{CPI10} gives the raw 10-year expectation
and column \texttt{CPI1} gives the realized inflation term used in the
adjustment. The adjusted output is written as a separate minimal table
\texttt{adjusted\_cpi10.csv} with columns
\texttt{(survey\_year, survey\_quarter, forecaster\_id, adjusted\_cpi10)}. This
adjusted object is the individual-level long-term inflation expectation
$x_{u,s}$ used in the forecast revision regression.

With quarterly inflation, the raw 10-year object is the average of 40 quarterly
inflation rates. For the first-quarter survey, there is no realized quarter
inside the 10-year window, so
\begin{equation}
\text{SPF-10YR-Q1}_t
\approx \frac{1}{40}
\sum_{j=0}^{39} \pi_{t-1+j \to t+j}.
\end{equation}
Define the long-term inflation expectation by survey date as
\begin{align}
x_{u,(y,1)} &\equiv \text{CPI10}_{u,(y,1)}, \\
x_{u,(y,2)} &\equiv
 \frac{40\cdot\text{CPI10}_{u,(y,2)} - \text{CPI1}_{u,(y,2)}}{39}, \\
x_{u,(y,3)} &\equiv
 \frac{40\cdot\text{CPI10}_{u,(y,3)}
 - \left(\text{CPI1}_{u,(y,2)} + \text{CPI1}_{u,(y,3)}\right)}{38}, \\
x_{u,(y,4)} &\equiv
 \frac{40\cdot\text{CPI10}_{u,(y,4)}
 - \left(\text{CPI1}_{u,(y,2)} + \text{CPI1}_{u,(y,3)}
 + \text{CPI1}_{u,(y,4)}\right)}{37}.
\end{align}
The corresponding forward-looking measurement equations are
\begin{align}
x_{u,(y,1)} &= \frac{1}{40}\sum_{j=0}^{39} \pi_{t-1+j \to t+j}
 + \varepsilon_{u,(y,1)}, \\
x_{u,(y,2)} &= \frac{1}{39}\sum_{j=0}^{38} \pi_{t-1+j \to t+j}
 + \varepsilon_{u,(y,2)}, \\
x_{u,(y,3)} &= \frac{1}{38}\sum_{j=0}^{37} \pi_{t-1+j \to t+j}
 + \varepsilon_{u,(y,3)}, \\
x_{u,(y,4)} &= \frac{1}{37}\sum_{j=0}^{36} \pi_{t-1+j \to t+j}
 + \varepsilon_{u,(y,4)}.
\end{align}
Under this quarterly representation, CPI1 enters the raw 10-year forecast with
weight $1/40$ once it is realized, but the adjusted object re-averages only over
the remaining forward-looking quarters. Thus Q2 rescales by $39$, Q3 rescales by
$38$, and Q4 rescales by $37$. The raw CPI10 measure therefore mechanically
contains more realized inflation in Q4 than in Q1.

\subsection{Construct inflation news}

For each forecast unit and survey date, inflation news is defined as the
realized inflation for the current quarter minus the nowcast of that same
quarter made in the previous survey. In the cleaned wide table
\texttt{forecast\_individual.csv}, the realized term is \texttt{CPI1} from the
current survey row, while the nowcast term is \texttt{CPI2} from the previous
survey row for the same unit. The constructed output is written as
\texttt{inflation\_news.csv} with columns
\texttt{(survey\_year, survey\_quarter, forecaster\_id, inflation\_news)}.
Thus,
\begin{equation}
n_{u,s} \equiv \text{CPI1}_{u,s} - \text{CPI2}_{u,s^{-}}.
\end{equation}
For individual forecasts, this means matching the current row
$(\textit{survey\_year}, \textit{survey\_quarter}, \textit{forecaster\_id})$ to
the lagged survey row for the same \texttt{forecaster\_id}, with the Q1 rollover
encoded in $s^{-}$. This notation makes the timing explicit: $n_{u,s}$ is always
dated at the current survey $s$, while the forecast component comes from the
previous survey $s^{-}$.

\subsection{Construct reputation measure}

A dedicated regression config file determines how the long-term inflation
expectation $x_{u,s}$ is defined for the forecast revision pipeline. The config
contains an ordered list of specifications to run, namely raw
\texttt{CPI10} and adjusted \texttt{CPI10}. If one specification selects raw
\texttt{CPI10}, then
\begin{equation}
x_{u,s} \equiv \text{CPI10}_{u,s}.
\end{equation}
If another specification selects adjusted \texttt{CPI10}, then $x_{u,s}$
equals the adjusted long-term inflation expectation constructed in the previous
subsection and stored in \texttt{adjusted\_cpi10.csv}. In either case, $u$
corresponds to \texttt{forecaster\_id}. Each configured definition of $x_{u,s}$
is propagated through its own versions of \texttt{reputation\_measure.csv} and
\texttt{regression\_dataset.csv}. Define the reputation measure $\rho$ from
\begin{equation}
x_{u,s} = \rho_{u,s}\left((1-q)\pi_{\text{target}} + q z_a\right)
 + (1-\rho_{u,s})\left((1-q)\pi_{\text{NE}} + q z_{\alpha}\right).
\end{equation}
Given values for $q$, $\pi_{\text{target}}$, $\pi_{\text{NE}}$, $z_a$, and
$z_{\alpha}$, this equation implies
\begin{equation}
\rho_{u,s} =
\frac{x_{u,s} - \left((1-q)\pi_{\text{NE}} + q z_{\alpha}\right)}
{\left((1-q)\pi_{\text{target}} + q z_a\right)
 - \left((1-q)\pi_{\text{NE}} + q z_{\alpha}\right)}.
\end{equation}
The implementation must load the parameter values from a config file rather
than hardcoding them in the construction procedure. It must also load the
configured list of long-term inflation expectation definitions from the
dedicated regression config file. For each observation where $x_{u,s}$ is
available, compute $\rho_{u,s}$ using the configured parameter values and the
selected definition of $x_{u,s}$ dated at the same survey $s$. The constructed
output for each specification is written as
\texttt{reputation\_measure.csv} with columns
\texttt{(survey\_year, survey\_quarter, forecaster\_id, rho)} in its own
specification-specific output directory.

\subsection{Construct regression dataset}

The regression dataset is indexed at the survey-date level. For each survey
date $s$, define the individual long-term inflation forecast revision for unit
$u$ using the same config-selected definition of $x_{u,s}$ by
\begin{equation}
r_{u,s} \equiv x_{u,s} - x_{u,s^{-}}.
\end{equation}
The regression dependent variable is the matched-sample mean of these
individual revisions across units observed in both $s$ and $s^{-}$. Let
$\mathcal{U}_s$ denote the set of units $u$ such that both $x_{u,s}$ and
$x_{u,s^{-}}$ are observed. Then the dependent variable is
\begin{equation}
\bar{r}_s \equiv \frac{1}{|\mathcal{U}_s|}\sum_{u \in \mathcal{U}_s} r_{u,s}.
\end{equation}
The same matched sample $\mathcal{U}_s$ is used to construct the survey-level
regressors. Define
\begin{align}
\bar{n}_s &\equiv \frac{1}{|\mathcal{U}_s|}\sum_{u \in \mathcal{U}_s} n_{u,s}, \\
\bar{\rho}_{s^{-}} &\equiv \frac{1}{|\mathcal{U}_s|}\sum_{u \in \mathcal{U}_s} \rho_{u,s^{-}}, \\
z_{2,s} &\equiv \bar{n}_s \bar{\rho}_{s^{-}}\left(1-\bar{\rho}_{s^{-}}\right), \\
z_{3,s} &\equiv \bar{n}_s \bar{\rho}_{s^{-}}\left(1-\bar{\rho}_{s^{-}}\right)^2, \\
z_{P,s} &\equiv \bar{n}_s \bar{\rho}_{s^{-}}^2\left(1-\bar{\rho}_{s^{-}}\right).
\end{align}
Thus the regression-ready dataset at survey date $s$ contains the dependent
variable $\bar{r}_s$ and the four survey-level regressors
$\bar{n}_s$, $z_{2,s}$, $z_{3,s}$, and $z_{P,s}$. The reputation term entering the
nonlinear regressors is the matched-sample mean previous-survey reputation
$\bar{\rho}_{s^{-}}$, computed using the same set of units $\mathcal{U}_s$.
The constructed survey-level dataset is written as
\texttt{regression\_dataset.csv}.

\section{OLS regressions for forecast revision}

The four OLS regressions are estimated at the survey level with a constant:
\begin{align}
\bar{r}_s &= \alpha_1 + \beta_1 \bar{n}_s + \varepsilon_{1,s}, \\
\bar{r}_s &= \alpha_2 + \beta_2 z_{2,s} + \varepsilon_{2,s}, \\
\bar{r}_s &= \alpha_3 + \beta_3 z_{3,s} + \varepsilon_{3,s}, \\
\bar{r}_s &= \alpha_P + \beta_P z_{P,s} + \varepsilon_{P,s}.
\end{align}
The estimation sample is restricted by the dedicated regression config file,
which specifies the start survey date and end survey date to use for both
estimation and fitted-value comparisons for each configured specification.

\subsection{Reported regression statistics}

For each of the four regressions, report a main table containing the slope
coefficient estimate, slope p value, adjusted $R^2$, RMSE, and sample size.
Report a second table directly below the main table containing intercept
inference for each model: intercept estimate, intercept standard error, and
intercept p value. Both tables should have one row per model and use the same
regression-ready sample used to estimate that model.

\subsection{Comparison figure}

Construct a figure showing the cumulative change in long-term inflation
forecast over the same config-selected sample window used in estimation. The
figure should plot the cumulative sum of the survey-level data series
$\bar{r}_s$ against the cumulative sums of the fitted counterparts produced by
the four regression models, using the same bounded survey window in each
series. Plot the data using a black dashed line.
Plot the fitted counterpart from Model 1 using a magenta solid line, Model 2
using a blue solid line, Model 3 using a red solid line, and placebo Model P
using a green solid line.

\section{Regression results}

The regression implementation writes results tables and a comparison figure to
\texttt{data-pipeline/output/forecast\_revision/} (worktree-relative). The table is
\texttt{regression\_results.csv}; it contains one row for each model and the
reported statistics described above: intercept estimate, intercept standard
error, intercept p value, slope estimate, slope p value, adjusted $R^2$, RMSE,
and sample size. The LaTeX output
\texttt{regression\_results.tex} contains two stacked tables: the main
forecast revision regression table and an intercept-inference table directly
under it. The figure is
\texttt{cumulative\_forecast\_revision\_comparison.png}; it plots the cumulative
change in long-term inflation forecast over the config-selected sample window
for the data series and
for the fitted counterparts from Models 1, 2, 3, and P using the specified line
styles and colors.

\section{Algorithms: download automation}

\begin{algorithm}[H]
\caption{\textsc{VariablePageURL}$(\textit{var})$}
\begin{algorithmic}[1]
\State \textbf{Require} $\textit{var}$: variable name as on Philadelphia Fed data-files index (e.g., NGDP, CPI10).
\State $\textit{path} \gets \text{lowercase}(\textit{var})$.
\If{$\textit{var} = \texttt{CPI}$}
    \State $\textit{path} \gets \texttt{cpi-spf}$. \Comment{exception per Philadelphia Fed URL}
\EndIf
\State \Return \texttt{https://www.philadelphiafed.org/surveys-and-data/} $\oplus$ $\textit{path}$.
\end{algorithmic}
\end{algorithm}

\begin{algorithm}[H]
\caption{\textsc{GetDownloadLinks}$(\textit{page\_url}, \textit{file\_types})$}
\begin{algorithmic}[1]
\State \textbf{Require} $\textit{file\_types}$: subset of
\Statex \quad $\{\texttt{dispersion}, \texttt{median\_level}, \texttt{mean\_level}, \texttt{median\_growth}, \texttt{mean\_growth}, \texttt{individual}\}$, and optionally \texttt{documentation}.
\State Fetch HTML at $\textit{page\_url}$.
\State Parse page for download links: match \texttt{<a>} tags allowing
\Statex \quad \texttt{href} to appear after other attributes (e.g.\ \texttt{class}); match link text or href to each requested type (e.g., ``Median Responses'' $\to$ \texttt{Median\_VAR\_Level.xlsx}, ``Measures of Cross-Sectional Forecast Dispersion'' $\to$ \texttt{Dispersion\_VAR.xlsx}).
\State Build list $L$ of $(\textit{filename}, \textit{download\_url})$ for each matched link; extract $\textit{filename}$ from the URL or link (stable name, e.g., \texttt{Median\_NGDP\_Level.xlsx}).
\State \Return $L$.
\end{algorithmic}
\end{algorithm}

\begin{algorithm}[H]
\caption{\textsc{DownloadOneFile}$(\textit{download\_url}, \textit{outpath}, \texttt{overwrite})$}
\begin{algorithmic}[1]
\State \textbf{Require} $\texttt{overwrite} \in \{\texttt{overwrite}, \texttt{skip-if-exists}\}$.
\If{$\texttt{overwrite} = \texttt{skip-if-exists}$ and file exists at $\textit{outpath}$}
    \State \Return (skip).
\EndIf
\State Download from $\textit{download\_url}$ to $\textit{outpath}$ (overwrite if file exists when $\texttt{overwrite} = \texttt{overwrite}$).
\State \Return $\textit{outpath}$.
\end{algorithmic}
\end{algorithm}

\begin{algorithm}[H]
\caption{\textsc{DownloadByVariableNames}$(\textit{variable\_names}, \textit{file\_types}, \textit{out\_dir}, \texttt{overwrite})$}
\begin{algorithmic}[1]
\State \textbf{Require} $\textit{out\_dir} = \texttt{data-pipeline/input/}$ (fixed by project layout).
\State \textbf{Require} $\textit{file\_types}$: which file types to download per variable (e.g., median level, mean level, dispersion, individual).
\State $\textit{written} \gets \emptyset$. \Comment{optional: track filenames already written to avoid duplicate writes}
\For{each $v$ in $\textit{variable\_names}$}
    \State $\textit{page\_url} \gets \textsc{VariablePageURL}(v)$.
    \State $L \gets \textsc{GetDownloadLinks}(\textit{page\_url}, \textit{file\_types})$.
    \For{each $(\textit{filename}, \textit{url})$ in $L$}
        \State $\textit{outpath} \gets \textit{out\_dir}$ joined with $\textit{filename}$.
        \State \textsc{DownloadOneFile}$(\textit{url}, \textit{outpath}, \texttt{overwrite})$.
        \State Add $\textit{outpath}$ to $\textit{written}$ (if not already present).
    \EndFor
\EndFor
\State \Return $\textit{written}$ (list of paths in \texttt{data-pipeline/input/}).
\end{algorithmic}
\end{algorithm}

\section{Algorithm: construct cleaned datasets (individual only)}

The following procedure summarizes the construction of SPF \emph{individual}
forecast datasets from raw Excel files to cleaned CSV. Raw files
\texttt{Individual\_VAR.xlsx} are assumed present in \texttt{data-pipeline/input/}
(e.g., produced by \textsc{DownloadByVariableNames} with \texttt{individual} in
\texttt{file\_types}). Output is CSV; \texttt{forecast\_individual} is written in
wide form (one row per survey/forecaster, one column per horizon).
Helper functions used by the main procedure are specified below.

\begin{algorithm}[H]
\caption{\textsc{LoadIndividualSheet}$(\textit{path})$}
\begin{algorithmic}[1]
\State \textbf{Require} $\textit{path}$: path to one \texttt{Individual\_VAR.xlsx} file.
\State Load first sheet: row 1 = header; first four columns YEAR, QUARTER, ID, INDUSTRY; remaining columns = horizon names.
\State Validate: at least YEAR, QUARTER, ID present; at least one horizon column.
\State Melt to long form: id\_vars = (YEAR, QUARTER, ID, INDUSTRY), value\_vars = horizon columns; create columns \texttt{horizon}, \texttt{value}.
\State Rename YEAR $\to$ survey\_year, QUARTER $\to$ survey\_quarter, ID $\to$ forecaster\_id.
\State \Return long-form table (survey\_year, survey\_quarter, forecaster\_id, industry if present, horizon, value).
\end{algorithmic}
\end{algorithm}

\begin{algorithm}[H]
\caption{\textsc{HorizonSortKey}$(\textit{col})$}
\begin{algorithmic}[1]
\State \textbf{Require} $\textit{col}$: horizon column name (e.g., CPI1, CPI10).
\State \textbf{Ensure} 10-year horizon (name ending with \texttt{10}) is ordered last.
\State \Return pair $(\textit{col}\,\text{ends with}\,\texttt{10}, \textit{col})$ for sorting (False before True, then by name).
\end{algorithmic}
\end{algorithm}

\begin{algorithm}[H]
\caption{\textsc{BuildForecastIndividual}$(\textit{df\_long})$}
\begin{algorithmic}[1]
\State \textbf{Require} $\textit{df\_long}$: concatenated long-form table from \textsc{LoadIndividualSheet}.
\State Pivot: index = (survey\_year, survey\_quarter, forecaster\_id); columns = horizon; values = value.
\State If the same horizon appears in more than one raw file, keep the first non-missing value.
\State Order horizon columns: sort by \textsc{HorizonSortKey} so 10-year column is last.
\State Final columns: (survey\_year, survey\_quarter, forecaster\_id, horizon\_1, horizon\_2, \ldots, horizon\_10).
\State \Return wide table (key columns as integers where applicable).
\end{algorithmic}
\end{algorithm}

\begin{algorithm}[H]
\caption{\textsc{BuildForecasterSurvey}$(\textit{df\_long})$}
\begin{algorithmic}[1]
\State \textbf{Require} $\textit{df\_long}$: concatenated long-form table (must contain survey\_year, survey\_quarter, forecaster\_id, INDUSTRY or industry).
\State Select distinct (survey\_year, survey\_quarter, forecaster\_id, industry); sort by (survey\_year, survey\_quarter, forecaster\_id).
\State \Return forecaster\_survey table; primary key (survey\_year, survey\_quarter, forecaster\_id).
\end{algorithmic}
\end{algorithm}

\begin{algorithm}[H]
\caption{\textsc{CleanIndividualTo3NF}$(\textit{input\_dir}, \textit{cleaned\_dir})$}
\begin{algorithmic}[1]
\State \textbf{Require} $\textit{input\_dir}$: path to input folder; $\textit{cleaned\_dir}$: path to cleaned output folder.
\State Find all \texttt{Individual\_*.xlsx} in $\textit{input\_dir}$.
\For{each file $\textit{path}$ in that list}
    \State $\textit{df} \gets \textsc{LoadIndividualSheet}(\textit{path})$.
    \State Append $\textit{df}$ to collection of long-form tables.
\EndFor
\State $\textit{df\_long} \gets$ concatenate all long-form tables.
\State $\textit{forecast\_individual} \gets \textsc{BuildForecastIndividual}(\textit{df\_long})$.
\State $\textit{forecaster\_survey} \gets \textsc{BuildForecasterSurvey}(\textit{df\_long})$.
\State Write $\textit{forecast\_individual}$ to $\textit{cleaned\_dir}$/\texttt{forecast\_individual.csv}, $\textit{forecaster\_survey}$ to $\textit{cleaned\_dir}$/\texttt{forecaster\_survey.csv}.
\State \Return the two tables (forecast\_individual, forecaster\_survey).
\end{algorithmic}
\end{algorithm}

\begin{algorithm}[H]
\caption{\textsc{GetQuarterSpecificValue}$(\textit{key}, \textit{horizon})$}
\begin{algorithmic}[1]
\State \textbf{Require} $\textit{key}$:
\Statex \quad $(\textit{forecast\_individual}, \textit{survey\_year}, \textit{survey\_quarter}, \textit{forecaster\_id})$.
\State \textbf{Require} $\textit{forecast\_individual}$: cleaned wide table.
\State Find row with key
\Statex \quad $(\textit{survey\_year}, \textit{survey\_quarter}, \textit{forecaster\_id})$.
\State \Return value in column $\textit{horizon}$ (e.g., \texttt{CPI1} or \texttt{CPI10}).
\end{algorithmic}
\end{algorithm}

\section{Algorithms: prepare dataset for forecast revision regression}

\subsection{Construct long-term inflation expectation}

\begin{algorithm}[H]
\caption{\textsc{AdjustCPI10Forecasts}$(\textit{forecast\_individual})$}
\begin{algorithmic}[1]
\State \textbf{Require} $\textit{forecast\_individual}$: cleaned wide table
\Statex \quad containing columns \texttt{CPI1} and \texttt{CPI10}.
\State Initialize output table for adjusted 10-year CPI forecasts.
\For{each row key $(\textit{survey\_year}, \textit{survey\_quarter}, \textit{forecaster\_id})$}
    \State $\textit{key} \gets (\textit{forecast\_individual}, \textit{survey\_year}, \textit{survey\_quarter}, \textit{forecaster\_id})$
    \State $\textit{cpi10} \gets \textsc{GetQuarterSpecificValue}(\textit{key}, \texttt{CPI10})$
    \If{$\textit{survey\_quarter} = 1$}
        \State $\textit{adjusted} \gets \textit{cpi10}$
    \ElsIf{$\textit{survey\_quarter} = 2$}
        \State $\textit{key\_q2} \gets (\textit{forecast\_individual}, \textit{survey\_year}, 2, \textit{forecaster\_id})$
        \State $\textit{cpi1\_q2} \gets \textsc{GetQuarterSpecificValue}(\textit{key\_q2}, \texttt{CPI1})$
        \State $\textit{adjusted} \gets (40\cdot\textit{cpi10} - \textit{cpi1\_q2})/39$
    \ElsIf{$\textit{survey\_quarter} = 3$}
        \State $\textit{key\_q2} \gets (\textit{forecast\_individual}, \textit{survey\_year}, 2, \textit{forecaster\_id})$
        \State $\textit{key\_q3} \gets (\textit{forecast\_individual}, \textit{survey\_year}, 3, \textit{forecaster\_id})$
        \State $\textit{cpi1\_q2} \gets \textsc{GetQuarterSpecificValue}(\textit{key\_q2}, \texttt{CPI1})$
        \State $\textit{cpi1\_q3} \gets \textsc{GetQuarterSpecificValue}(\textit{key\_q3}, \texttt{CPI1})$
        \State $\textit{adjusted} \gets (40\cdot\textit{cpi10} - (\textit{cpi1\_q2} + \textit{cpi1\_q3}))/38$
    \Else
        \State $\textit{key\_q2} \gets (\textit{forecast\_individual}, \textit{survey\_year}, 2, \textit{forecaster\_id})$
        \State $\textit{key\_q3} \gets (\textit{forecast\_individual}, \textit{survey\_year}, 3, \textit{forecaster\_id})$
        \State $\textit{key\_q4} \gets (\textit{forecast\_individual}, \textit{survey\_year}, 4, \textit{forecaster\_id})$
        \State $\textit{cpi1\_q2} \gets \textsc{GetQuarterSpecificValue}(\textit{key\_q2}, \texttt{CPI1})$
        \State $\textit{cpi1\_q3} \gets \textsc{GetQuarterSpecificValue}(\textit{key\_q3}, \texttt{CPI1})$
        \State $\textit{cpi1\_q4} \gets \textsc{GetQuarterSpecificValue}(\textit{key\_q4}, \texttt{CPI1})$
        \State $\textit{adjusted} \gets (40\cdot\textit{cpi10} - (\textit{cpi1\_q2} + \textit{cpi1\_q3} + \textit{cpi1\_q4}))/37$
    \EndIf
    \State Write $\textit{adjusted}$ to the output table for that row key.
\EndFor
\State \Return table of adjusted 10-year CPI forecasts.
\end{algorithmic}
\end{algorithm}

\subsection{Construct inflation news}

\begin{algorithm}[H]
\caption{\textsc{ConstructInflationNews}$(\textit{forecast\_individual})$}
\begin{algorithmic}[1]
\State \textbf{Require} $\textit{forecast\_individual}$: cleaned wide table
\Statex \quad containing columns \texttt{CPI1} and \texttt{CPI2}.
\State Initialize output table for inflation news.
\For{each row key $(\textit{survey\_year}, \textit{survey\_quarter}, \textit{forecaster\_id})$}
    \State $\textit{current\_key} \gets (\textit{forecast\_individual}, \textit{survey\_year}, \textit{survey\_quarter}, \textit{forecaster\_id})$
    \State $\textit{cpi1\_current} \gets \textsc{GetQuarterSpecificValue}(\textit{current\_key}, \texttt{CPI1})$
    \If{$\textit{survey\_quarter} = 1$}
        \State $\textit{lagged\_key} \gets (\textit{forecast\_individual}, \textit{survey\_year} - 1, 4, \textit{forecaster\_id})$
    \Else
        \State $\textit{lagged\_key} \gets (\textit{forecast\_individual}, \textit{survey\_year}, \textit{survey\_quarter} - 1, \textit{forecaster\_id})$
    \EndIf
    \State $\textit{cpi2\_lagged} \gets \textsc{GetQuarterSpecificValue}(\textit{lagged\_key}, \texttt{CPI2})$
    \State $\textit{inflation\_news} \gets \textit{cpi1\_current} - \textit{cpi2\_lagged}$
    \State Write $\textit{inflation\_news}$ to the output table for that row key.
\EndFor
\State \Return table of inflation news.
\end{algorithmic}
\end{algorithm}

\begin{algorithm}[H]
\caption{\textsc{SelectLongTermInflationExpectation}$(\textit{forecast\_individual}, \textit{regression\_config})$}
\begin{algorithmic}[1]
\State \textbf{Require} $\textit{regression\_config}$ contains \texttt{x\_definitions}.
\For{each \texttt{x\_definition} in \texttt{x\_definitions}}
\If{\texttt{x\_definition} = \texttt{raw\_cpi10}}
    \State Set \texttt{x} equal to \texttt{CPI10} from \texttt{forecast\_individual}.
\ElsIf{\texttt{x\_definition} = \texttt{adjusted\_cpi10}}
    \State Set \texttt{x} equal to the adjusted long-term inflation expectation.
\Else
    \State Raise an error for unsupported \texttt{x\_definition}.
\EndIf
\State \Return table indexed by
\Statex \quad $(\textit{survey\_year}, \textit{survey\_quarter}, \textit{forecaster\_id}, \textit{x})$.
\EndFor
\end{algorithmic}
\end{algorithm}

\subsection{Construct reputation measure}

\begin{algorithm}[H]
\caption{\textsc{ConstructReputationMeasure}$(\textit{x\_table}, \textit{config})$}
\begin{algorithmic}[1]
\State \textbf{Require} $\textit{x\_table}$: table containing config-selected long-term inflation expectation \texttt{x}.
\State \textbf{Require} $\textit{x\_table}$ has columns
\Statex \quad $(\textit{survey\_year}, \textit{survey\_quarter}, \textit{forecaster\_id}, \textit{x})$.
\State \textbf{Require} $\textit{config}$: parameter source containing
\Statex \quad \texttt{q}, \texttt{pi\_target}, \texttt{pi\_NE}, \texttt{z\_a}, \texttt{z\_alpha}.
\State Load parameter values from $\textit{config}$.
\State $\textit{target\_term} \gets (1-\texttt{q})\cdot\texttt{pi\_target} + \texttt{q}\cdot\texttt{z\_a}$.
\State $\textit{ne\_term} \gets (1-\texttt{q})\cdot\texttt{pi\_NE} + \texttt{q}\cdot\texttt{z\_alpha}$.
\State $\textit{denominator} \gets \textit{target\_term} - \textit{ne\_term}$.
\State Initialize output table for reputation measure.
\For{each row in $\textit{x\_table}$}
    \State $\textit{x} \gets$ value in column \texttt{x}.
    \State $\rho \gets (\textit{x} - \textit{ne\_term}) / \textit{denominator}$.
    \State Write $\rho$ to the output table with the same keys as the input row.
\EndFor
\State \Return table containing the input keys and \texttt{rho}.
\end{algorithmic}
\end{algorithm}

\subsection{Construct regression dataset}

\begin{algorithm}[H]
\caption{\textsc{ConstructRegressionDataset}$(\textit{x\_table}, \textit{inflation\_news}, \textit{reputation\_measure})$}
\begin{algorithmic}[1]
\State \textbf{Require} $\textit{x\_table}$: table containing config-selected long-term inflation expectation \texttt{x}.
\State \textbf{Require} $\textit{inflation\_news}$: table containing \texttt{inflation\_news}.
\State \textbf{Require} $\textit{reputation\_measure}$: table containing \texttt{rho}.
\State \textbf{Require} each table is indexed by
\Statex \quad $(\textit{survey\_year}, \textit{survey\_quarter}, u)$ with a common unit index $u$.
\State Initialize output table for the regression dataset.
\For{each survey date $s=(y,q)$ in sorted survey order}
    \State Determine previous survey date $s^{-}$.
    \State Build matched sample $\mathcal{U}_s$ of units $u$ such that
    \Statex \quad $x_{u,s}$, $x_{u,s^{-}}$, $n_{u,s}$, and $\rho_{u,s^{-}}$ are all observed.
    \If{$|\mathcal{U}_s| = 0$}
        \State Continue to next survey date.
    \EndIf
    \For{each $u \in \mathcal{U}_s$}
        \State $r_{u,s} \gets x_{u,s} - x_{u,s^{-}}$.
    \EndFor
    \State $\bar{r}_s \gets \frac{1}{|\mathcal{U}_s|}\sum_{u \in \mathcal{U}_s} r_{u,s}$.
    \State $\bar{n}_s \gets \frac{1}{|\mathcal{U}_s|}\sum_{u \in \mathcal{U}_s} n_{u,s}$.
    \State $\bar{\rho}_{s^{-}} \gets \frac{1}{|\mathcal{U}_s|}\sum_{u \in \mathcal{U}_s} \rho_{u,s^{-}}$.
    \State $z_{2,s} \gets \bar{n}_s \bar{\rho}_{s^{-}}\left(1-\bar{\rho}_{s^{-}}\right)$.
    \State $z_{3,s} \gets \bar{n}_s \bar{\rho}_{s^{-}}\left(1-\bar{\rho}_{s^{-}}\right)^2$.
    \State $z_{P,s} \gets \bar{n}_s \bar{\rho}_{s^{-}}^2\left(1-\bar{\rho}_{s^{-}}\right)$.
    \State Write $(s, \bar{r}_s, \bar{n}_s, \bar{\rho}_{s^{-}}, z_{2,s}, z_{3,s}, z_{P,s}, |\mathcal{U}_s|)$
    \Statex \quad to the regression dataset.
\EndFor
\State \Return survey-level regression dataset.
\end{algorithmic}
\end{algorithm}

\section{Algorithms: OLS regressions for forecast revision}

\subsection{Reported regression statistics}

\begin{algorithm}[H]
\caption{\textsc{RunForecastRevisionRegressions}$(\textit{regression\_dataset}, \textit{regression\_config})$}
\begin{algorithmic}[1]
\State \textbf{Require} $\textit{regression\_dataset}$: survey-level table containing
\Statex \quad $(\textit{survey\_year}, \textit{survey\_quarter}, \bar{r}_s, \bar{n}_s, z_{2,s}, z_{3,s}, z_{P,s})$.
\State \textbf{Require} $\textit{regression\_config}$ contains the sample start and end survey dates.
\State Keep only survey dates inside the configured sample window for the current specification.
\State Initialize output table for regression statistics.
\State Estimate Model 1 by OLS with a constant:
\Statex \quad $\bar{r}_s = \alpha_1 + \beta_1 \bar{n}_s + \varepsilon_{1,s}$.
\State Estimate Model 2 by OLS with a constant:
\Statex \quad $\bar{r}_s = \alpha_2 + \beta_2 z_{2,s} + \varepsilon_{2,s}$.
\State Estimate Model 3 by OLS with a constant:
\Statex \quad $\bar{r}_s = \alpha_3 + \beta_3 z_{3,s} + \varepsilon_{3,s}$.
\State Estimate placebo Model P by OLS with a constant:
\Statex \quad $\bar{r}_s = \alpha_P + \beta_P z_{P,s} + \varepsilon_{P,s}$.
\For{each estimated model $m \in \{1,2,3,P\}$}
    \State Extract intercept estimate, intercept standard error, intercept p value.
    \State Extract slope estimate, slope p value, adjusted $R^2$, RMSE, and sample size.
    \State Write one row of statistics for model $m$ to the output table.
\EndFor
\State \Return regression-statistics table and fitted values from the four models.
\end{algorithmic}
\end{algorithm}

\subsection{Comparison figure}

\begin{algorithm}[H]
\caption{\textsc{PlotCumulativeForecastRevisionComparison}$(\textit{regression\_dataset}, \textit{fitted\_values}, \textit{regression\_config})$}
\begin{algorithmic}[1]
\State \textbf{Require} fitted values from Models 1, 2, 3, and P for each survey date in
\Statex \quad $\textit{regression\_dataset}$.
\State \textbf{Require} $\textit{regression\_config}$ contains the sample start and end survey dates.
\State Keep survey dates inside the configured sample window; sort by survey date.
\State Compute cumulative data series as the running sum of $\bar{r}_s$.
\State Compute cumulative fitted series for each model as the running sum of its fitted values.
\State Construct figure with five lines:
\Statex \quad Data = black dashed line.
\Statex \quad Model 1 fitted counterpart = magenta solid line.
\Statex \quad Model 2 fitted counterpart = blue solid line.
\Statex \quad Model 3 fitted counterpart = red solid line.
\Statex \quad Placebo Model P fitted counterpart = green solid line.
\State Label horizontal axis with survey date and vertical axis with cumulative change
\Statex \quad in long-term inflation forecast over the configured sample window.
\State \Return comparison figure.
\end{algorithmic}
\end{algorithm}

\section{Results and comparison}

The forecast revision pipeline is run twice using the two
\texttt{x\_definitions} listed in
\texttt{data-pipeline/config/forecast\_revision.json}: raw \texttt{CPI10} and
adjusted \texttt{CPI10}. This produces two specification-specific regression
tables, two specification-specific regression figures, and a set of
survey-level matched-sample comparison figures. The embedded objects below are
generated by the pipeline and written under
\texttt{data-pipeline/output/forecast\_revision/}.

\subsection{Regression results: raw CPI10}

\input{\forecastrevisionoutputroot/raw_cpi10/regression_results.tex}
The intercept inference table is reported directly under the main regression
table in the same rendered block above.

\begin{figure}[h]
\centering
\includegraphics[width=\textwidth]{\forecastrevisionoutputroot/raw_cpi10/cumulative_forecast_revision_comparison.png}
\caption{Cumulative forecast revision comparison using raw CPI10.}
\end{figure}

\subsection{Regression results: adjusted CPI10}

\input{\forecastrevisionoutputroot/adjusted_cpi10/regression_results.tex}
The intercept inference table is reported directly under the main regression
table in the same rendered block above.

\begin{figure}[h]
\centering
\includegraphics[width=\textwidth]{\forecastrevisionoutputroot/adjusted_cpi10/cumulative_forecast_revision_comparison.png}
\caption{Cumulative forecast revision comparison using adjusted CPI10.}
\end{figure}

\subsection{Survey-level matched-sample comparison figures}

\begin{figure}[H]
\centering
\includegraphics[width=\textwidth]{\forecastrevisionoutputroot/comparison/long_term_forecast_revision_comparison.png}
\caption{Survey-level matched-sample mean long-term forecast revision: raw versus adjusted CPI10.}
\end{figure}

\begin{figure}[H]
\centering
\includegraphics[width=\textwidth]{\forecastrevisionoutputroot/comparison/reputation_measure_comparison.png}
\caption{Survey-level matched-sample mean previous-survey reputation: raw versus adjusted CPI10.}
\end{figure}

\begin{figure}[H]
\centering
\includegraphics[width=\textwidth]{\forecastrevisionoutputroot/comparison/regressor_1_comparison.png}
\caption{Survey-level matched-sample mean regressor 1: raw versus adjusted CPI10.}
\end{figure}

\begin{figure}[H]
\centering
\includegraphics[width=\textwidth]{\forecastrevisionoutputroot/comparison/regressor_2_comparison.png}
\caption{Survey-level matched-sample mean regressor 2: raw versus adjusted CPI10.}
\end{figure}

\begin{figure}[H]
\centering
\includegraphics[width=\textwidth]{\forecastrevisionoutputroot/comparison/regressor_3_comparison.png}
\caption{Survey-level matched-sample mean regressor 3: raw versus adjusted CPI10.}
\end{figure}

\begin{figure}[H]
\centering
\includegraphics[width=\textwidth]{\forecastrevisionoutputroot/comparison/regressor_P_comparison.png}
\caption{Survey-level matched-sample mean placebo regressor: raw versus adjusted CPI10.}
\end{figure}

\section{Notes}

Before running the pipeline, synchronize dependencies from
\texttt{data-pipeline/pyproject.toml} using \texttt{uv sync}. Run all pipeline
and test commands with \texttt{uv run ...} so execution uses the managed
environment.

Obtaining official Excel files may be implemented by \textsc{DownloadByVariableNames}: given variable names (e.g., CPI10, NGDP), the procedure uses \textsc{VariablePageURL} to form each variable's data page URL (e.g., \url{https://www.philadelphiafed.org/surveys-and-data/cpi10}), \textsc{GetDownloadLinks} to parse the page for Individual (and optionally other) links, and \textsc{DownloadOneFile} to save each file to \texttt{data-pipeline/input/} under the same filename as on the site. The currently implemented cleaned individual-level outputs are \texttt{forecast\_individual}, \texttt{forecaster\_survey}, \texttt{adjusted\_cpi10}, \texttt{inflation\_news}, and specification-specific versions of \texttt{reputation\_measure} and \texttt{regression\_dataset}, all written as CSV. The dedicated regression config file \texttt{data-pipeline/config/forecast\_revision.json} specifies both the sample window for the regression analysis, the ordered list of long-term inflation expectation definitions to compare, and a worktree-relative output root for forecast revision artifacts. Under the forecast revision regression design, the analysis results are written to \texttt{data-pipeline/output/forecast\_revision/}, including specification-specific regression tables and figures plus the survey-level matched-sample comparison figures. Level/mean/median and probability tables remain out of scope unless added explicitly later.
Codebook tables and figures are generated in \texttt{data-pipeline/report/codebook.qmd}; they are not part of the construction pseudocode.

\end{document}
